% Labels of supplementary_SR.tex, exported from supplementary_SR.aux by
% export_xr_labels.py. Read by preamble_SR.tex only when the .aux
% itself is absent (Overleaf drops .aux files from uploaded ZIPs).
% Do not edit: regenerate with .\build_SR.ps1.
\newlabel{fig:signal_tree}{{S1}{2}{\textbf {Where the mutation signal lives, gene by gene: the three regimes of \cref {eq:pid} as a decision tree over the four training conditions.} Each read is directional: it follows the sign of the deciding contrast, and none is established, because no fusion contrast excludes zero or falls inside the $\pm 0.02$ equivalence margin (Table~2). Regime~(i), encoder saturation (\textit {TP53}, \textit {EGFR}): the embedding already carries what the vocabulary names. Regime~(ii), complementarity: its additive sub-form (ii-a) is empty in LUAD, and its non-additive sub-form (ii-b), tested on \textit {KRAS} by the conditional co-activation contrast, returns empty (\cref {tab:kras_coactivation}). Regime~(iii), out-of-ontology (\textit {STK11}, \textit {KEAP1}, \textit {RBM10}): the vocabulary converts into no gain, for three different reasons that one Inconclusive label conflates---label-space mismatch, modality ceiling and collapse to prior. The cohort size each read would need for confirmation is in \cref {tab:design_power}}{figure.caption.2}{supplementary_SR.pdf}}
\newlabel{fig:signal_tree@cref}{{[suppfigure][1][]S1}{[1][2][]2}{}{}{supplementary_SR.pdf}}
\newlabel{fig:kras_lepidic_solid_proxy}{{S2}{3}{\textbf {The pattern pair the Choquet head selects for KRAS, and the textbook morphology it stands for.} Left, the three components of a KRAS-driven invasive mucinous adenocarcinoma slide (a tumour-cell layer growing along alveoli, dense sheet-like masses, extracellular mucin). Right, what the $6$-d ANORAK vector retains of it: lepidic and solid light up together, and the mucin has no pattern. The positive interaction $I(\text {lepidic}\!\times \!\text {solid})\!=\!{+}0.032$ the $2$-additive measure learns on \textit {KRAS} is the aggregator paying for that co-occurrence, the histological proxy for invasive mucinous adenocarcinoma, the KRAS-enriched subtype.\cite {rekhtman2013krassolid} The index is a parameter of the fitted measure shared by every slide; whether it buys accuracy where the pair co-activates is the conditional test of \cref {tab:kras_coactivation}, which returns empty at this cohort size. The full Choquet read-out is \cref {fig:choquet_kras}}{figure.caption.3}{supplementary_SR.pdf}}
\newlabel{fig:kras_lepidic_solid_proxy@cref}{{[suppfigure][2][]S2}{[1][2][]3}{}{}{supplementary_SR.pdf}}
\newlabel{fig:ui_stack}{{S3}{4}{\textbf {The released per-slide tooling (\lchai {}~v2 interface)}, for KRAS on TCGA-55-7815, a case whose cohort verdict is Inconclusive yet whose slide-level prediction is $P(\mathrm {mut})=0.965$. The top strip carries the population verdict and the slide-level prediction, with the molecular-testing flag raised whenever $\auroc <0.70$. The six read-outs are materialised from a single forward pass: \textbf {L1} ABMIL attention, \textbf {L2} pattern overlay, \textbf {L3} SHAP embedding/pattern split, \textbf {L4} Choquet--Shapley singleton values and pairwise interaction indices, \textbf {L5} channel ablation, and \textbf {L6} the fuzzy-label language output with its knowledge-graph-grounded references. The tooling is released as software and is not evaluated in this study (Methods). All imagery is from the public TCGA-LUAD cohort; the system is a research tool and is not for clinical diagnosis. The cohort \auroc {} printed in the strip ($0.609$) is the earlier single-split estimate served by the deployment at the time of capture (Supplementary Note~2) and differs from the nested cross-validation values of Table~1}{figure.caption.4}{supplementary_SR.pdf}}
\newlabel{fig:ui_stack@cref}{{[suppfigure][3][]S3}{[1][2][]4}{}{}{supplementary_SR.pdf}}
\newlabel{fig:pipeline_glioma}{{S4}{5}{\textbf {The glioma MRI\,$\rightarrow $\,genotype pipeline as a one-to-one transfer of the LUAD architecture.} An unconstrained learned encoder $\mathbf {E}$---a 2.5D CNN with \abmil {} over axial MRI slices, occupying the position \ctranspath {} with tile-level \abmil {} holds in LUAD---and the expert source vocabulary $\mathbf {P}$---the VASARI semantic MR features, expert-scored on the TCGA-LGG arm and propagated from the annotated seed for the BraTS MGMT arm, occupying the position ANORAK growth patterns hold in LUAD---feed the same four conditions: B2 encoder only, B3 ontology only, PI concat fusion and FC $2$-additive Choquet. Measured outcomes are in \cref {tab:crossdomain_merged} (MGMT, full power) and \cref {tab:glioma_domainb} (TCGA-LGG pilot)}{figure.caption.5}{supplementary_SR.pdf}}
\newlabel{fig:pipeline_glioma@cref}{{[suppfigure][4][]S4}{[1][2][]5}{}{}{supplementary_SR.pdf}}
\newlabel{fig:pipeline_ecg}{{S5}{6}{\textbf {The electrocardiogram\,$\rightarrow $\,physiological-state pipeline running the identical four-condition protocol in a different modality family.} The encoder $\mathbf {E}$---a 1D convolutional network over the raw twelve-lead signal---and the source vocabulary $\mathbf {P}$---the $71$ graded SCP-ECG statements---feed B2, B3, PI and FC exactly as \ctranspath {} embeddings and ANORAK patterns do in LUAD. Both targets are Conclusive and encoder-saturated: the fusion conditions match the encoder to within a fold's standard deviation while the statement channel alone trails}{figure.caption.6}{supplementary_SR.pdf}}
\newlabel{fig:pipeline_ecg@cref}{{[suppfigure][5][]S5}{[1][2][]6}{}{}{supplementary_SR.pdf}}
\newlabel{fig:two_channels}{{S6}{6}{\textbf {The two-channel architecture in full.} Tile embeddings from the frozen \ctranspath {} encoder and $6$-d fuzzy pattern probabilities from the ANORAK-trained classifier remain disjoint until the classifier head, which is what makes the decomposition of \cref {eq:pid} well posed as a statement about two sources}{figure.caption.7}{supplementary_SR.pdf}}
\newlabel{fig:two_channels@cref}{{[suppfigure][6][]S6}{[1][2][]6}{}{}{supplementary_SR.pdf}}
\newlabel{fig:choquet_kras}{{S7}{7}{\textbf {Full Choquet read-out for KRAS, rendered on slide TCGA-55-7815.} The six singleton values $\phi _k$ and all $15$ pairwise interaction indices $I_{jk}$ learned by the $2$-additive fuzzy measure, from which the lepidic\,$\times $\,solid interaction of $+0.032$ reported in the main text is read. These are parameters of the fitted measure and are shared by every slide the model scores; the figure is titled by slide because that is the interface the read-out is served through, not because the values are estimated per slide. The conditional test of whether this interaction buys accuracy where the pair co-activates is \cref {tab:kras_coactivation}, and it is unresolved}{figure.caption.8}{supplementary_SR.pdf}}
\newlabel{fig:choquet_kras@cref}{{[suppfigure][7][]S7}{[1][2][]7}{}{}{supplementary_SR.pdf}}
\newlabel{fig:glioma_example}{{S8}{7}{\textbf {A real preoperative multiparametric MRI from a TCGA-LGG glioma} (IDH-mutant, 1p/19q-codeleted)---the four co-registered sequences (T1, T1-Gd, T2, FLAIR) the encoder $\mathbf {E}$ ingests, at the tumour-centroid axial slice with the expert segmentation contour (red). Below, the neuroradiologist VASARI scores for the same case, the source channel $\mathbf {P}$: the radiological counterpart of the ANORAK growth patterns in LUAD}{figure.caption.9}{supplementary_SR.pdf}}
\newlabel{fig:glioma_example@cref}{{[suppfigure][8][]S8}{[1][2][]7}{}{}{supplementary_SR.pdf}}
\newlabel{fig:ecg_ptbxl}{{S9}{8}{\textbf {A real twelve-lead electrocardiogram from PTB-XL} (a $37$-year-old woman, sinus rhythm, reported normal), in the standard clinical $3\times 4$ format with a lead-II rhythm strip ($25$\,mm/s, $10$\,mm/mV). This is the raw $100$\,Hz signal the encoder ingests. Age and sex are recorded for every case but absent from the $71$ SCP-ECG statements forming the source channel, which is what makes the recovery reported in \cref {tab:crossdomain_merged} leakage-free.\cite {wagner2020ptbxl}}{figure.caption.10}{supplementary_SR.pdf}}
\newlabel{fig:ecg_ptbxl@cref}{{[suppfigure][9][]S9}{[1][2][]8}{}{}{supplementary_SR.pdf}}
\newlabel{fig:pid_estimator}{{S10}{8}{\textbf {Why the atoms are read as directional signatures rather than estimated numerically.} The spread of the per-fold single-channel \auroc {}s from which any plug-in atom estimate would be built, against the size of the pattern-side effect such an estimate would have to resolve}{figure.caption.11}{supplementary_SR.pdf}}
\newlabel{fig:pid_estimator@cref}{{[suppfigure][10][]S10}{[1][2][]8}{}{}{supplementary_SR.pdf}}
\newlabel{tab:levels_of_analysis}{{S1}{9}{\textbf {The four levels of analysis, the quantities estimated at each, and the scope of the claims each supports.} The Conclusive/Inconclusive verdict and the three-regime taxonomy are population-level objects; the per-slide read-outs (L1--L6, Methods) are slide-level and tile-level objects; neither transfers to the other. KRAS is the canonical illustration, and in both directions: a confident per-slide call does not lift the population verdict, and a population-level interaction weight does not descend into per-case evidence (\cref {tab:kras_coactivation}). Only the name of the unit and whether it is a bag or a single instance change when the framework moves to glioma MRI or to the electrocardiogram, which is the precise sense in which the framework, and not the model weights, transfers}{table.caption.13}{supplementary_SR.pdf}}
\newlabel{tab:levels_of_analysis@cref}{{[supptable][1][]S1}{[1][9][]9}{}{}{supplementary_SR.pdf}}
\newlabel{tab:literature_comparison}{{S2}{10}{\textbf {Per-gene \auroc {} against the published mutation-prediction competition on LUAD.} This work reports six driver genes simultaneously on the public TCGA-LUAD cohort under a single MIL aggregator; prior work typically reports two to three genes on either the same public cohort or on proprietary cohorts an order of magnitude larger. Within \textbf {This work}, bold marks the better of the two aggregators per gene. No value is bolded \emph {across} studies: the protocols are not comparable, mixing held-out test sets with cross-validation and public with proprietary cohorts, so a column maximum would not identify a better method. A per-gene \auroc {} also does not answer the question this paper asks---whichever study reports the higher one, the number does not say whether the signal lies inside the source ontology or outside it, and that is what decides whether the next experiment is a larger cohort, a richer vocabulary or a different modality. No study below reports a regime or a per-target verdict. The comparison is developed in Supplementary Note~2}{table.caption.14}{supplementary_SR.pdf}}
\newlabel{tab:literature_comparison@cref}{{[supptable][2][]S2}{[1][9][]10}{}{}{supplementary_SR.pdf}}
\newlabel{tab:cptac_transfer}{{S3}{11}{\textbf {Cohort transfer to CPTAC-LUAD: the frozen TCGA-trained pipeline applied off-the-shelf} ($498$ slides, $106$ genotyped patients). Read this way the TCGA verdicts do not transfer. What fails, however, is the transported decision function and not the tissue, so no assignment is licensed as near-intrinsic on this evidence, but neither is the coupling itself cohort-local. Per-gene values, the aggregator and stain-normalisation controls and the power caveat are set out below the table}{table.caption.15}{supplementary_SR.pdf}}
\newlabel{tab:cptac_transfer@cref}{{[supptable][3][]S3}{[1][9][]11}{}{}{supplementary_SR.pdf}}
\newlabel{tab:cptac_refit}{{S4}{12}{\textbf {Re-fitting the four conditions inside CPTAC-LUAD, with only the representation carried over from TCGA.} This is the comparison the frozen-pipeline transfer of \cref {tab:cptac_transfer} cannot make, and it separates the two readings that transfer alone leaves open. The recovered values, the field-of-view control and the regime adjudication under the noise-aware rule are set out below the table}{table.caption.16}{supplementary_SR.pdf}}
\newlabel{tab:cptac_refit@cref}{{[supptable][4][]S4}{[1][9][]12}{}{}{supplementary_SR.pdf}}
\newlabel{tab:atom_operationalization}{{S5}{12}{\textbf {How each atom of \cref {eq:pid} is read from the model.} Every atom is mapped to the model-side signature that evidences it, the directional decision rule on the four training conditions, the interpretability level(s) that expose it, and the regime it drives. The signatures are read, not numerically estimated; atom names follow the Bertschinger convention.\cite {bertschinger2014unique}}{table.caption.17}{supplementary_SR.pdf}}
\newlabel{tab:atom_operationalization@cref}{{[supptable][5][]S5}{[1][9][]12}{}{}{supplementary_SR.pdf}}
\newlabel{tab:kras_coactivation}{{S6}{13}{\textbf {The conditional test of regime~(ii-b) on KRAS.} The Choquet head is credited with recovering a lepidic$\,\times \,$solid synergy, which predicts that its advantage over additive fusion should concentrate where the two patterns co-occur. Cutting the cohort by a model-independent measure of that co-occurrence does not reproduce the prediction: the contrast is flat across strata and the high${-}$low difference disagrees in sign between the two protocols. The stratifying variable is itself associated with KRAS, so what fails is the aggregator step, not the morphological premise}{table.caption.18}{supplementary_SR.pdf}}
\newlabel{tab:kras_coactivation@cref}{{[supptable][6][]S6}{[1][9][]13}{}{}{supplementary_SR.pdf}}
\newlabel{tab:design_power}{{S7}{14}{\textbf {What cohort size a conditional regime claim requires.} Prospective design calculation for the contrast of \cref {tab:kras_coactivation}: target effects are stipulated a priori and the variance is taken from the hierarchical bootstrap at $n\!=\!444$, so this is design guidance rather than post-hoc power on the observed null. At TCGA-LUAD scale only implausibly large effects are resolvable, which is why we report the KRAS conditional result as a bounded null and why we recommend that interaction-level claims be pre-registered against a size target of this kind}{table.caption.19}{supplementary_SR.pdf}}
\newlabel{tab:design_power@cref}{{[supptable][7][]S7}{[1][9][]14}{}{}{supplementary_SR.pdf}}
\newlabel{tab:literature_power_audit}{{S8}{15}{\textbf {Reported aggregator gains against the resolution of the designs that report them.} The power argument of \cref {tab:design_power} is not specific to this cohort. Applying the same calculation to published mutation-prediction results, using each study's own AUROC and class counts, places every audited gain below the minimum effect its design could detect, and none of those studies reports an interval or a test. Two independent lines from within the field agree: where comparisons are formally tested most are not significant, and the same aggregator varies more between reimplementations than any claimed improvement}{table.caption.20}{supplementary_SR.pdf}}
\newlabel{tab:literature_power_audit@cref}{{[supptable][8][]S8}{[1][9][]15}{}{}{supplementary_SR.pdf}}
\newlabel{tab:controls}{{S9}{16}{\textbf {Control grid: tumour versus normal, permuted labels, patient sex and \textit {TP53}, on three seeds under both protocols.} The negative control lands on chance everywhere and the positive control clears it on both channels, so the low absolute values of \cref {tab:mutation_luad_resultssec} are a property of the cohort and encoder rather than of a defect in the pipeline. \textit {TP53} reproduces the five-seed main arm to within $0.011$, which we report as replication rather than as proof of the pipeline, a burden the tumour arm now carries instead. The sex control neither passes nor fails: the design cannot resolve the effect it was chosen to detect, so it is reported as underpowered---the verdict \cref {tab:literature_power_audit} issues on the published literature, applied here to ourselves}{table.caption.21}{supplementary_SR.pdf}}
\newlabel{tab:controls@cref}{{[supptable][9][]S9}{[1][9][]16}{}{}{supplementary_SR.pdf}}
\newlabel{tab:mutation_prevalence}{{S10}{17}{Mutation prevalence in the analysable TCGA-LUAD cohort ($n\!=\!458$ tumour slides from $444$ patients, derived in Methods from the $687$-image inventory). The Conclusive-verdict column anticipates the cohort-level outcome of the cross-validation reported in Results}{table.caption.22}{supplementary_SR.pdf}}
\newlabel{tab:mutation_prevalence@cref}{{[supptable][10][]S10}{[1][9][]17}{}{}{supplementary_SR.pdf}}
\newlabel{tab:viz_folds_methods}{{S11}{17}{Conclusive/Inconclusive verdict bands and their assignment to the six target genes under the five-fold patient-stratified cross-validation protocol used throughout}{table.caption.23}{supplementary_SR.pdf}}
\newlabel{tab:viz_folds_methods@cref}{{[supptable][11][]S11}{[1][9][]17}{}{}{supplementary_SR.pdf}}
\newlabel{tab:threshold_sensitivity}{{S12}{18}{\textbf {Sensitivity of every categorical claim to the two free thresholds.} The Conclusive cut of \cref {eq:conclusive_verdict} and the equivalence margin that gates regimes~(i) and~(iii) are the only two parameters in the framework that are set rather than estimated. Both are swept here. No LUAD target is Conclusive at any cut at or above $0.65$ except \textit {TP53} on one protocol, and no regime read changes over $\delta \!\in \!\{.01,.02,.03\}$. At this cohort size the binding constraint is interval width rather than the margin: the median cell would need $\sim \!14{,}000$ patients before equivalence could be certified at $\pm .02$}{table.caption.24}{supplementary_SR.pdf}}
\newlabel{tab:threshold_sensitivity@cref}{{[supptable][12][]S12}{[1][9][]18}{}{}{supplementary_SR.pdf}}
\newlabel{tab:cptac_power}{{S13}{19}{Observed positive-set sizes for the six target genes in the CPTAC-LUAD cohort-transfer test, from the non-silent mutation frequencies in cBioPortal \texttt {luad\_cptac\_2020} ($n\!=\!110$). These fix, in advance of the imaging read-out, which targets the transfer can estimate: in the genotyped-with-imaging subset ($n\!=\!106$) \textit {TP53} and \textit {EGFR} clear the $n_+\!=\!32$ floor and \textit {KRAS} sits at it}{table.caption.25}{supplementary_SR.pdf}}
\newlabel{tab:cptac_power@cref}{{[supptable][13][]S13}{[1][9][]19}{}{}{supplementary_SR.pdf}}
\newlabel{tab:pid_estimator}{{S14}{19}{\textbf {The plug-in atom estimator the analysis specifies but declines to use.} For each term of \cref {eq:pid}, what it captures, how a plug-in estimator would compute it, and the role it plays in the regime assignment. No numerical atom value is reported anywhere in this work}{table.caption.26}{supplementary_SR.pdf}}
\newlabel{tab:pid_estimator@cref}{{[supptable][14][]S14}{[1][9][]19}{}{}{supplementary_SR.pdf}}
\newlabel{tab:crossdomain_protocol}{{S15}{20}{Cohorts, source ontologies, encoders and pre-registered expected regime assignments for the two additional modalities, fixed before the read-outs were computed}{table.caption.27}{supplementary_SR.pdf}}
\newlabel{tab:crossdomain_protocol@cref}{{[supptable][15][]S15}{[1][9][]20}{}{}{supplementary_SR.pdf}}
\newlabel{tab:levels_crossdomain}{{S16}{21}{The four levels of analysis instantiated in each of the three modalities, showing that only the name of the unit and whether it is a bag or a single instance change}{table.caption.28}{supplementary_SR.pdf}}
\newlabel{tab:levels_crossdomain@cref}{{[supptable][16][]S16}{[1][9][]21}{}{}{supplementary_SR.pdf}}
\newlabel{tab:glioma_domainb}{{S17}{21}{\textbf {Glioma MRI, full per-condition read-out.} The $n\!\sim \!63$ TCGA-LGG pilot with expert-scored VASARI (IDH, 1p/19q and a patient-sex control) together with the full-power MGMT test on BraTS~2021 ($n\!=\!577$). Only the MGMT row is a main-text result (panel~(a) of \cref {tab:crossdomain_merged}); the pilot rows are directional reads at a cohort size that cannot adjudicate a regime (Supplementary Note~1), and no regime is asserted from them}{table.caption.29}{supplementary_SR.pdf}}
\newlabel{tab:glioma_domainb@cref}{{[supptable][17][]S17}{[1][9][]21}{}{}{supplementary_SR.pdf}}
\newlabel{tab:pid_regimes}{{S18}{22}{\textbf {The three regimes as atoms of \cref {eq:pid}, with the model-side signature that evidences each atom in LUAD.} Every regime in the table is the pre-registered expectation against which the adjudication of Table~2 was run, not an established assignment: no fusion contrast excludes zero and none establishes equivalence at the $\pm 0.02$ margin, so all six genes remain unresolved. The table is retained because the \emph {implied remedy} column is what the decomposition adds over a cohort-mean \auroc {}, and because it separates the three states one Inconclusive label conflates (empty, unestimable, unconverted)}{table.caption.30}{supplementary_SR.pdf}}
\newlabel{tab:pid_regimes@cref}{{[supptable][18][]S18}{[1][9][]22}{}{}{supplementary_SR.pdf}}
\newlabel{tab:ecg_domainc}{{S19}{23}{Full per-condition electrocardiogram read-out with the directional-criterion discussion, condensed into panel~(b) of \cref {tab:crossdomain_merged}}{table.caption.31}{supplementary_SR.pdf}}
\newlabel{tab:ecg_domainc@cref}{{[supptable][19][]S19}{[1][9][]23}{}{}{supplementary_SR.pdf}}
